\chapter{Introduction}
\label{ch:introduction}

% ─────────────────────────────────────────────────────────────────────────────
\section{What This Manual Covers}
\label{sec:intro-covers}
% ─────────────────────────────────────────────────────────────────────────────

This manual is the authoritative reference for the \texttt{gert.home} Domain Kit---a
specialized vocabulary layer for authoring household maintenance runbooks within the gert
framework. The kit implements the \textbf{Property} model, providing step types, scheduling
primitives, and evidence-capture patterns designed for homeowners and property managers tracking
household maintenance workflows.

The manual covers:

\begin{itemize}
  \item The \texttt{Property} model and role taxonomy
  \item \texttt{gert.home} step types: \texttt{home.routine},
        \texttt{home.incident}, \texttt{home.delegate}
  \item Delegation and away-mode enforcement via \texttt{IsAwayModeActive()}
  \item Evidence types: \texttt{photo}, \texttt{note}, \texttt{checklist}, \texttt{none}
  \item Consumable tracking: replacement schedules and stock alerts
  \item Complete field reference for all \texttt{gert.home} schema elements
\end{itemize}

This is \emph{not} a general introduction to gert or Domain Kits. For core gert concepts
(runbook structure, core step types, the Domain Kit model), see the \emph{gert Design
Document}. For guidance on developing your own Domain Kit, see the \emph{Domain Kit Development
Guide}.

% ─────────────────────────────────────────────────────────────────────────────
\section{Audience and Prerequisites}
\label{sec:intro-audience}
% ─────────────────────────────────────────────────────────────────────────────

This manual is written for:

\begin{description}
  \item[Homeowners] --- The primary audience. Homeowners manage their own property file,
        define zones and assets, schedule routines, and review evidence after each completed
        task.
  \item[Property managers] --- Professionals managing multiple properties. They author one
        \texttt{.home.yaml} file per property, configure delegations for caretakers, and
        review compliance bundles.
  \item[Household contractors] --- Gardeners, plumbers, and other tradespeople who are
        assigned routines via \texttt{executor.role} hints or active delegation blocks. They
        execute tasks and submit evidence but do not modify the property file.
\end{description}

\paragraph{Prerequisites.}
You should be familiar with:

\begin{itemize}
  \item \textbf{gert core concepts} --- runbook structure, core step types (\texttt{cli},
        \texttt{manual}, \texttt{tool}, \texttt{branch}, \texttt{iterate}), variables, and the
        execution model. See \emph{gert Design Document}, §1--§3.
  \item \textbf{Domain Kit concepts} --- what a Domain Kit is, how Kit-specific step types
        compile (lower) into core primitives, and the relationship between authoring vocabulary
        and runtime execution. See \emph{gert Design Document}, §4.
  \item \textbf{YAML syntax} --- \texttt{gert.home} property files are authored in YAML. You
        do not need to know JSON Schema or Go.
\end{itemize}

% ─────────────────────────────────────────────────────────────────────────────
\section{What is a Property?}
\label{sec:intro-concept}
% ─────────────────────────────────────────────────────────────────────────────

\textbf{Property} is the top-level household unit in \texttt{gert.home}. It is the single
source of truth for everything that must be maintained: the named areas (Zones), tracked
physical items (Assets), recurring tasks (Routines), reactive repair workflows
(IncidentTemplates), items on replacement schedules (Consumables), and temporary assignees
(Delegations).

\paragraph{Origin.}
The Property model is drawn from household management systems and facilities management
software. The core insight is that maintenance accountability requires an explicit owner and a
structured record---not just a to-do list. gert.home brings that structure into the runbook
execution engine.

\paragraph{Why the Property Model Matters.}
Without the Property model, household maintenance workflows face:
\begin{itemize}
  \item \textbf{Reactive chaos} --- tasks are remembered only when something breaks, not on
        a schedule.
  \item \textbf{No audit trail} --- there is no record of who did what, when, or what
        evidence was captured.
  \item \textbf{Missed delegations} --- when the homeowner is away, there is no explicit
        assignment of responsibility to a delegate.
\end{itemize}

The Property model resolves this by making household accountability explicit and
machine-readable through the runbook schema.

\paragraph{The Property in \texttt{gert.home}.}
The \texttt{gert.home} Domain Kit implements the Property model through:
\begin{itemize}
  \item A structured \texttt{property} block: \texttt{name}, \texttt{zones[]}
  \item Tracked items: \texttt{assets[]}, \texttt{consumables[]}
  \item Scheduled work: \texttt{routines[]} with cadence and evidence requirements
  \item Reactive repair: \texttt{incident\_templates[]} with sequential and parallel steps
  \item Away-mode: \texttt{delegations[]} with scoped assignments and permissions
\end{itemize}

% ─────────────────────────────────────────────────────────────────────────────
\section{How \texttt{gert.home} Implements the Property Model}
\label{sec:intro-implementation}
% ─────────────────────────────────────────────────────────────────────────────

The \texttt{gert.home} kit extends gert core with domain-specific step types and
metadata fields. These compile down to core primitives at build time so the gert runtime
never sees \texttt{gert.home}-specific vocabulary.

\subsection{Step Types}

\begin{description}
  \item[\texttt{home.routine}] --- A scheduled recurring maintenance task. Lowers to an
        \texttt{iterate} step containing a \texttt{manual} step. Supports fixed-interval
        (\texttt{every: 7d}) and seasonal cadences.
  \item[\texttt{home.incident}] --- A reactive repair workflow with sequential, parallel, and
        branching steps. Each incident step lowers to a \texttt{manual}, \texttt{branch}, or
        \texttt{parallel} core step as appropriate.
  \item[\texttt{home.delegate}] --- An away-mode delegation block. Checked at run-time via
        \texttt{IsAwayModeActive()}; when active, substitutes the delegate identity as executor
        for the assigned routines or zones.
\end{description}

\subsection{Metadata Fields}

\texttt{gert.home} property files include additional top-level blocks:

\begin{description}
  \item[\texttt{property}] --- Display name and zone list for the household unit.
  \item[\texttt{assets}] --- Tracked physical items with installation date and warranty info.
  \item[\texttt{consumables}] --- Items on a replacement schedule with stock tracking.
  \item[\texttt{delegations}] --- Away-mode assignment blocks with date ranges and permissions.
\end{description}

\subsection{Evidence Types}

\texttt{gert.home} defines four evidence types for routine and incident steps:

\begin{description}
  \item[\texttt{photo}] --- A photograph required to prove task completion (e.g., clear pool
        water, mowed lawn).
  \item[\texttt{note}] --- A free-text note from the executor (e.g., cost of repair,
        measurements recorded).
  \item[\texttt{checklist}] --- A structured checklist that must be completed item by item
        before the step closes.
  \item[\texttt{none}] --- No evidence required; the step is self-certifying.
\end{description}

\subsection{Delegation Integration}

When a \texttt{home.delegate} block is active, the \texttt{gert.home} runtime calls
\texttt{IsAwayModeActive()} before each routine execution. If the delegation is active and the
routine or zone is in the delegate's \texttt{assigns} list, the delegate's identity is
substituted as executor. When the delegation expires, the homeowner resumes accountability
automatically.

% ─────────────────────────────────────────────────────────────────────────────
\section{Document Structure}
\label{sec:intro-structure}
% ─────────────────────────────────────────────────────────────────────────────

The remaining chapters are organized as follows:

\begin{description}
  \item[Chapter 1: gert.home Concepts] --- The Property model, role taxonomy, the
        delegation chain, and cadence/SLA concepts.
  \item[Chapter 2: Schema Reference] --- The \texttt{gert.home} kit schema: step types,
        property metadata fields, and evidence types.
  \item[Chapter 3: Authoring Guide] --- How to write a \texttt{.home.yaml} property file
        from scratch, with a complete pool-maintenance example.
  \item[Chapter 4: gert Integration] --- Role-based gating, away-mode enforcement, cadence
        checking, and the audit trail.
  \item[Chapter 5: Evidence and Tracing] --- The evidence model, run directory structure,
        timeline projection, SHA-256 verification, and compliance bundles.
  \item[Chapter 6: Reference] --- Complete field reference tables for all
        \texttt{gert.home} schema elements and environment variables.
\end{description}
